\documentclass[../main.tex]{subfiles}
\begin{document}
\chapter{Time Integration and Constraint Drift}
\label{chap:integration-drift}

\section{State of the Art}
Constrained mechanical models can be integrated as differential-algebraic equations, reduced to independent coordinates, or converted into an ordinary differential equation using differentiated constraints. These choices interact with time-stepping and drift control but are not the same decision. Baumgarte's feedback method modifies the acceleration constraint, while projection methods correct a numerical state back toward the constraint manifold \cite{baumgarte1972stabilization,ascherpetzold1998}. Structure-preserving integration provides another important perspective for long-time simulation, though it does not by itself select a loop-constraint treatment \cite{hairer2006geometric}. This is a focused orientation based on these foundational sources and the SciPy 1.18.1 documentation accessed 2026-10-07, not an exhaustive review of current constrained-mechanics solvers.

\section{Separate the model from the stepper}
The ideal closed-chain equations are
\begin{align}
 M(q)\ddot q + J(q)^T\lambda &= Q(q,\dot q,t),\\
 \Phi(q)&=0.
\end{align}
They are a differential-algebraic system. Differentiating the constraint twice yields
\begin{equation}
 J(q)\ddot q = -\dot J(q,\dot q)\dot q.
\end{equation}
Solving this acceleration-level saddle-point system at each state evaluation determines \(\ddot q\) and \(\lambda\). Returning \([\dot q,\ddot q]\) then defines an ODE for the state \(y=[q,\dot q]^T\), but integration error can still move the state away from \(\Phi=0\) and \(J\dot q=0\). It is therefore important to report both the chosen differential formulation and the time integrator.

\section{Drift control}
\subsection{Baumgarte stabilization}
Let \(\Phi\) be the position residual. Baumgarte stabilization prescribes a damped second-order error equation,
\begin{equation}
 \ddot\Phi+2\zeta\omega_B\dot\Phi+\omega_B^2\Phi=0,
\end{equation}
where \(\zeta>0\) and \(\omega_B>0\) are numerical gains. Using \(\dot\Phi=J\dot q\) gives the stabilized acceleration constraint
\begin{equation}
 J\ddot q=-\dot J\dot q-2\zeta\omega_B J\dot q-\omega_B^2\Phi.
\end{equation}
The last two terms are feedback corrections, not physical forces or part of the ideal model. Increasing \(\omega_B\) may improve constraint-error decay but also introduce fast time scales and stiffness. Gains must be evaluated with the integration method and residuals, not selected from a universal rule.

\subsection{Projection}
Position projection solves a nonlinear correction problem that returns \(q\) toward \(\Phi(q)=0\). A local least-norm correction satisfies
\begin{equation}
 J(q_k)\Delta q=-\Phi(q_k),\qquad q_{k+1}=q_k+\Delta q.
\end{equation}
Velocity projection at a corrected position similarly solves \(J\Delta\dot q=-J\dot q\). Projection may be used for initialization or after accepted integration steps. It changes the numerical state and can affect invariants; its frequency and norm should be reported. The notebook projects only the initial position and velocity, not every later state.

\subsection{Other formulations}
Coordinate reduction eliminates dependent coordinates and constraints locally, but requires a suitable parameterization and may become singular or cumbersome as the mechanism changes configuration. Penalty methods replace exact constraints with stiff forces proportional to constraint error; they avoid explicit multipliers but introduce penalty parameters and potentially severe stiffness. A DAE solver treats algebraic constraints as part of the integrated problem and has solver-specific consistency requirements. These are formulation choices; each can be paired with different time discretizations. DAE solver theory and stiff integration are treated by Ascher and Petzold and by Hairer and Wanner \cite{ascherpetzold1998,hairerwanner1996}.

\section{Time-stepper and tolerances}
For the demonstration, the acceleration solve with Baumgarte feedback defines an ODE integrated by SciPy \texttt{solve\_ivp}. Its default method is RK45, an explicit Runge--Kutta method of order 5(4). Relative and absolute tolerances control local error estimates, while \texttt{max\_step} bounds accepted step size; \texttt{t\_eval} specifies stored output times and does not force internal steps at those points \cite{scipy_solve_ivp}. Explicit Runge--Kutta methods are appropriate first choices for non-stiff problems; the documentation recommends implicit methods such as Radau or BDF when stiffness is present. A Baumgarte gain that creates fast correction dynamics may require a stiff solver or a different drift treatment.

For the five-link run, the Baumgarte settings are \(\zeta=1\) and \(\omega_B=20\,\mathrm{s}^{-1}\). The ODE solver uses relative tolerance \(10^{-8}\), absolute tolerance \(10^{-10}\), and maximum step \(0.01\,\mathrm{s}\) over ten seconds. These are reproducible example settings, not a claim of optimality. A numerical study should vary tolerances, maximum step, and stabilization gains and compare constraint errors and quantities of interest.

\section{Algorithmic outline}
\begin{enumerate}
 \item Project approximate initial positions onto \(\Phi(q)=0\), then project velocities onto \(J(q)\dot q=0\).
 \item At each state evaluation, compute \(\Phi\), \(J\), \(\dot J\dot q\), and the applied generalized forces.
 \item Solve the stabilized KKT system for acceleration and multipliers.
 \item Return the state derivative formed by the generalized velocity and the solved acceleration.
 \item Integrate with an adaptive solver and record status, tolerances, solver evaluations, and saved output times.
 \item Independently evaluate position and velocity constraint residuals along the trajectory; adjust the method only after inspecting those diagnostics.
\end{enumerate}

\subsection{Notebook execution sequence}

The canonical five-link notebook implements the above as an ODE right-hand side. Written without language-specific syntax, one evaluation proceeds as follows:
\begin{enumerate}
\item Split $y=[q^T,\dot q^T]^T$ into the 15 generalized positions and 15 velocities.
\item Evaluate $\Phi(q)$, $J(q)$, $\gamma(q,\dot q)=\dot J\dot q$, and check that $J$ has the expected rank 13. Stop with a diagnostic if the constraints lose rank.
\item Form the stabilized lower-block right-hand side
\begin{equation}
r_c=-\gamma-2\zeta\omega_BJ\dot q-\omega_B^2\Phi.
\label{eq:five-link-stabilized-rhs-ch9}
\end{equation}
\item Solve the block system with $M$, $J^T$, $J$, and a zero lower-right block for $\ddot q$ and $\lambda$, using $[Q_{\mathrm{total}}^T,r_c^T]^T$ as its right-hand side.
\item Return $\dot y=[\dot q^T,\ddot q^T]^T$. The multipliers are algebraic outputs, not additional integrated state variables.
\end{enumerate}

For the recorded baseline, the initial state is projected once and then integrated for $10\,\mathrm{s}$ with RK45, relative tolerance $10^{-8}$, absolute tolerance $10^{-10}$, and maximum internal step $0.01\,\mathrm{s}$. The solver is asked to save 2,001 evenly spaced output samples; those requested times do not force its internal adaptive steps. After the run, compute $\|\Phi(q_k)\|_\infty$, $\|J(q_k)\dot q_k\|_\infty$, and grounded-pose error at every saved state, and record solver status and RHS evaluation count. The notebook then repeats the initial-value problem with tighter tolerances and with $\omega_B$ reduced from $20$ to $10\,\mathrm{s}^{-1}$, comparing residual maxima and maximum coordinate difference from baseline. These repeats are sensitivity checks, not a formal convergence proof.

\section{Summary}
Constraint enforcement, drift control, and time integration are separate choices. This book's example uses an acceleration-level KKT solve, initial projection, Baumgarte feedback, and adaptive RK45. The next chapter turns these choices into a verification and validation plan.

\ifSubfilesClassLoaded{\backmatter\bookbibliography}{}
\end{document}